% ============================================================
%  KICH BAN VIDEO CA NHAN --- KHANH
%  Vai tro: Backend Architecture / API / Pipeline / Frontend + Ket luan toan bo de an
%  Files: app_api.py (1893L) | run_phone_cam.py (1314L)
%  Thoi luong muc tieu: 14-16 phut (có kết luận toàn đề án)
% ============================================================
\documentclass[12pt, a4paper]{article}
\usepackage[utf8]{inputenc}
\usepackage[T5]{fontenc}
\usepackage[vietnamese]{babel}
\usepackage[margin=2cm, top=2.5cm, bottom=2.5cm]{geometry}
\usepackage{parskip}
\usepackage[dvipsnames]{xcolor}
\usepackage{tcolorbox}
\tcbuselibrary{skins, breakable, listings}
\usepackage{listings}
\usepackage{booktabs}
\usepackage{tabularx}
\usepackage{fontawesome5}
\usepackage{amsmath}
\usepackage{pifont}
\newcommand{\cmark}{\ding{51}}

\lstset{
  basicstyle=\ttfamily\footnotesize,
  keywordstyle=\color{NavyBlue}\bfseries,
  commentstyle=\color{OliveGreen}\itshape,
  stringstyle=\color{BrickRed},
  numbers=left, numberstyle=\tiny\color{gray}, numbersep=6pt,
  frame=single, breaklines=true, showstringspaces=false,
  tabsize=4, backgroundcolor=\color{gray!7},
}

\usepackage{titlesec}
\titleformat{\section}[block]
  {\large\bfseries\color{MidnightBlue!80!black}}{\thesection.}{0.5em}{}[\titlerule]
\titleformat{\subsection}[block]
  {\normalsize\bfseries\color{MidnightBlue!70!black}}{\thesubsection}{0.5em}{}

\usepackage{fancyhdr}
\pagestyle{fancy}\fancyhf{}
\lhead{\textcolor{MidnightBlue!80!black}{\textbf{Kịch Bản Video --- KHÁNH (Backend / API / Kết Luận)}}}
\rhead{\textcolor{gray}{XLA Face Privacy System}}
\cfoot{\thepage}
\usepackage[hidelinks]{hyperref}

%% ---- Environments ----
\newtcolorbox{say}[1][]{
  title={\faMicrophone\ \textbf{NÓI}},
  colback=MidnightBlue!5, colframe=MidnightBlue!60,
  fonttitle=\small\bfseries, left=4pt, right=4pt, breakable, #1}

\newtcolorbox{doact}[1][]{
  title={\faMousePointer\ \textbf{THAO TÁC}},
  colback=OliveGreen!5, colframe=OliveGreen!60,
  fonttitle=\small\bfseries, left=4pt, right=4pt, breakable, #1}

\newtcolorbox{cam}[1][]{
  title={\faVideo\ \textbf{CAMERA}},
  colback=Goldenrod!8, colframe=Goldenrod!80,
  fonttitle=\small\bfseries, left=4pt, right=4pt, breakable, #1}

\newtcolorbox{code}[2][]{
  title={\faCode\ \textbf{CODE --- #2}},
  colback=MidnightBlue!4, colframe=MidnightBlue!45,
  fonttitle=\small\bfseries\color{MidnightBlue}, left=4pt, right=4pt, breakable, #1}

\newtcolorbox{warn}[1][]{
  title={\faExclamationTriangle\ \textbf{CHÚ Ý}},
  colback=red!4, colframe=red!60,
  fonttitle=\small\bfseries, left=4pt, right=4pt, breakable, #1}

\newtcolorbox{timing}[1]{
  colback=MidnightBlue!8, colframe=MidnightBlue!60,
  title={\faClock\ \textbf{#1}},
  fonttitle=\small\bfseries\color{MidnightBlue!80!black}, left=4pt, right=4pt}

\newcommand{\fpath}[1]{\texttt{\textcolor{purple}{#1}}}
\newcommand{\lnum}[1]{\colorbox{yellow!40}{\texttt{\small dòng~#1}}}
\newcommand{\pause}{\textcolor{gray}{\textit{[dừng 1 giây]}}}
\newcommand{\camlook}{\textcolor{Goldenrod}{\textbf{[nhìn camera]}}}
\newcommand{\screen}{\textcolor{OliveGreen}{\textbf{[nhìn màn hình]}}}

% ============================================================
\begin{document}
% ============================================================

\begin{titlepage}
\centering\vspace*{1.5cm}
{\Huge\bfseries\color{MidnightBlue!80!black} KỊCH BẢN VIDEO\par}
\vspace{0.3cm}
{\LARGE\bfseries\color{MidnightBlue!70!black} KHÁNH --- Backend API + Kết Luận Toàn Đề Án\par}
\vspace{0.5cm}
{\large\color{gray} XLA Face Privacy System --- Đồ Án Cuối Kỳ\par}
\vspace{1.2cm}

\begin{tcolorbox}[colback=MidnightBlue!6, colframe=MidnightBlue!80, width=0.9\textwidth]
\centering\large
\textbf{Em trực tiếp viết 2 files + đảm nhiệm phần kết luận toàn đề án:}\\[0.4cm]
\renewcommand{\arraystretch}{1.5}
\begin{tabular}{llr}
  \hline
  \fpath{src/app\_api.py} & FastAPI Server: 40+ endpoints, Pipeline, Streaming & \textbf{1893 dòng} \\
  \fpath{src/run\_phone\_cam.py} & FrameReader + Phone Camera Integration & \textbf{1314 dòng} \\
  \hline
  \multicolumn{2}{l}{\textbf{Tổng cộng}} & \textbf{3207 dòng} \\
  \hline
\end{tabular}\\[0.4cm]
\textit{Đây là codebase lớn nhất trong nhóm.
Khánh cũng là người trình bày cuối cùng --- kết luận toàn bộ đề án.}
\end{tcolorbox}

\vspace{0.8cm}
\begin{tcolorbox}[colback=gray!6, colframe=gray!50, width=0.9\textwidth]
\textbf{Thời lượng mục tiêu:} 14--16 phút\\[0.2cm]
\textbf{Cần chuẩn bị:}
Server \texttt{uvicorn} chạy cổng 8000 \textbullet{}
VS Code mở \fpath{app\_api.py} + \fpath{run\_phone\_cam.py} \textbullet{}
Browser mở \texttt{dashboard.html} + \texttt{face-registry.html} + \texttt{secure-archive.html} \textbullet{}
Swagger UI: \texttt{localhost:8000/docs} \textbullet{}
Camera kết nối
\end{tcolorbox}

\vspace{0.8cm}
\begin{warn}
\textbf{Lưu ý cho Khánh:}
\fpath{app\_api.py} 1893 dòng --- không thể show hết. Chiến lược: show 5--6 điểm quan trọng và giải thích tại sao chúng quan trọng.
Phần kết luận toàn đề án: chuẩn bị số liệu cụ thể (749+759+3207 dòng = 4725 dòng, etc.)
\end{warn}

\end{titlepage}

\tableofcontents\newpage

% ============================================================
\section{CHUẨN BỊ TRƯỚC KHI BẤM REC}
% ============================================================

\begin{doact}
\textbf{Bước 1 --- Khởi động server:}
\begin{lstlisting}[language=bash, numbers=none]
cd C:\Users\itent\Desktop\xla_demo
.venv\Scripts\activate
uvicorn src.app_api:app --host 0.0.0.0 --port 8000 --reload
\end{lstlisting}
Chờ thấy: \texttt{Application startup complete} trên terminal.
\end{doact}

\begin{doact}
\textbf{Bước 2 --- VS Code:}
Tab 1: \fpath{src/app\_api.py} (1893 dòng)\\
Tab 2: \fpath{src/run\_phone\_cam.py} (1314 dòng)\\
Font size 16--18.
\end{doact}

\begin{doact}
\textbf{Bước 3 --- Browser (3 tab):}
Tab 1: \texttt{http://localhost:8000/admin/dashboard.html}\\
Tab 2: \texttt{http://localhost:8000/admin/face-registry.html}\\
Tab 3: \texttt{http://localhost:8000/admin/secure-archive.html}\\
Tab 4 (để Q\&A): \texttt{http://localhost:8000/docs} (Swagger auto-generated)
\end{doact}

\begin{cam}
Camera ngang tầm mắt. Ánh sáng từ phía trước. Nói rõ ràng, xương sống thẳng.
Video này kết thúc toàn bộ đề án --- tone tự tin, không vội vàng.
\end{cam}

\newpage
% ============================================================
\section{0:00--0:15 --- MỞ ĐẦU}
% ============================================================

\begin{timing}{0:00 -- 0:15}
Nhìn camera, mở đầu.
\end{timing}

\begin{say}
``Xin chào thầy cô. Em là Khánh.

Em phụ trách phần \textbf{Backend API và hệ thống Pipeline xử lý video real-time}.
Video này em trình bày 2 files em viết,
sau đó kết luận toàn bộ đề án của nhóm.''
\end{say}

% ============================================================
\section{0:15--1:30 --- GIỚI THIỆU ĐÓNG GÓP}
% ============================================================

\begin{doact}
Hover chuột qua từng tab file khi nói.
\end{doact}

\begin{say}
``Em viết 2 files lớn nhất trong đề án:

\textbf{app\_api.py} --- 1893 dòng.
Đây là trung tâm của toàn hệ thống --- FastAPI server.
40+ HTTP endpoints, WebSocket, MJPEG stream, pipeline orchestration.
Mọi thứ Đạt và Khoa viết đều được gọi thông qua server này.

\textbf{run\_phone\_cam.py} --- 1314 dòng.
Trong file này có class \texttt{FrameReader} --- thread đọc camera non-blocking.
Cũng có pipeline xử lý khi chạy từ command line.

Tổng cộng 3207 dòng --- nhiều nhất nhóm.

Em sẽ chỉ rõ 6 điểm quan trọng nhất:
Admin authentication, PipelineSession class, FrameReader,
main processing loop, streaming, và CRUD endpoints.''
\end{say}

% ============================================================
\section{1:30--3:00 --- APP\_API.PY --- ADMIN AUTHENTICATION}
% ============================================================

\begin{timing}{1:30 -- 3:00}
Chuyển sang \fpath{src/app\_api.py}.
\end{timing}

\begin{doact}
Click tab \fpath{app\_api.py}.
Nhấn \texttt{Ctrl+G}, gõ \textbf{394}, Enter.
\end{doact}

\begin{say}
``Điểm đầu tiên: \textbf{Admin Authentication flow}.

\fpath{app\_api.py} không tự xử lý password --- nó delegate hết cho \fpath{admin\_auth.py} của Đạt.
Nhưng em viết phần JWT token và middleware protect endpoints.''
\end{say}

\begin{code}{app\_api.py dòng 394--448 --- Admin auth: token + middleware}
\begin{lstlisting}[language=Python, firstnumber=394]
def _issue_admin_token() -> str:          # dong 394
    """Issue a short-lived admin JWT (no library)."""
    payload = {
        "sub": "admin",
        "iat": int(time.time()),
        "exp": int(time.time()) + 3600,   # 1 hour
    }
    token_b = json.dumps(payload).encode()
    return base64.urlsafe_b64encode(token_b).decode()

def _verify_admin_token(token: str) -> bool:   # dong 408
    try:
        payload = json.loads(base64.urlsafe_b64decode(token))
        return (payload.get("sub") == "admin" and
                time.time() < payload.get("exp", 0))
    except Exception:
        return False

async def _require_admin_auth(           # dong 426: middleware
    request: Request,
    x_admin_token: Optional[str] =
        Header(None, alias="X-Admin-Token"),
) -> None:
    if x_admin_token and _verify_admin_token(x_admin_token):
        return
    # Also check cookie
    cookie_tok = request.cookies.get("admin_token", "")
    if cookie_tok and _verify_admin_token(cookie_tok):
        return
    raise HTTPException(status_code=401,
                        detail="Admin authentication required")
\end{lstlisting}
\end{code}

\begin{say}
``\textbf{Dòng 394}: \texttt{\_issue\_admin\_token} tạo token JSON base64.
Token chứa: sub="admin", thời điểm tạo, thời điểm hết hạn (1 tiếng).

\textbf{Dòng 408}: \texttt{\_verify\_admin\_token} decode base64, parse JSON,
verify sub="admin" và chưa expired.

\textbf{Dòng 426}: \texttt{\_require\_admin\_auth} là dependency FastAPI.
Mọi endpoint cần quyền admin đều khai báo \texttt{Depends(\_require\_admin\_auth)}.
Kiểm tra header \texttt{X-Admin-Token} hoặc cookie \texttt{admin\_token}.
Nếu cả hai đều fail: HTTP 401.

Điểm quan trọng: chỉ sau khi Đạt's \texttt{verify\_admin} trong \fpath{admin\_auth.py}
trả True em mới gọi \texttt{\_issue\_admin\_token} ở dòng 394.
Không có password đúng $\to$ không có token $\to$ không qua middleware dòng 426.''
\end{say}

% ============================================================
\section{3:00--6:00 --- PIPELINESESSION CLASS --- LIVE VIDEO PIPELINE}
% ============================================================

\begin{timing}{3:00 -- 6:00}
Nhấn \texttt{Ctrl+G}, gõ \textbf{512}, Enter.
\end{timing}

\begin{say}
``Điểm thứ hai và quan trọng nhất: class \texttt{PipelineSession}.''
\end{say}

\subsection{3:00--4:30 --- PipelineSession \_\_init\_\_, start, stop, patch\_config}

\begin{code}{app\_api.py dòng 512--590 --- PipelineSession: pipeline lifecycle management}
\begin{lstlisting}[language=Python, firstnumber=512]
class PipelineSession:
    """
    Manages one live video pipeline: camera -> detect -> anonymize ->
    stream. One instance per active pipeline.
    """

    def __init__(self, camera_index: int,
                 preset_name: str = "balanced") -> None:  # dong 515
        self.camera_index = camera_index
        self.preset_name  = preset_name
        self.stop_event   = threading.Event()
        self.cfg_lock     = threading.Lock()
        self.live_cfg     = _LiveConfig(preset_name=preset_name)
        self._thread: Optional[threading.Thread] = None
        self.frame_buffer = _FrameBuffer()        # output buffer

    def start(self) -> None:                      # dong 571
        self.stop_event.clear()
        self._thread = threading.Thread(
            target=self._run, daemon=True)
        self._thread.start()

    def stop(self) -> None:                       # dong 575
        self.stop_event.set()                     # signal stop
        if self._thread:
            self._thread.join(timeout=10.0)       # wait graceful shutdown

    def is_alive(self) -> bool:                   # dong 580
        return (self._thread is not None and
                self._thread.is_alive())

    def patch_config(self,                        # dong 583
                     patch: Dict[str, Any]) -> None:
        with self.cfg_lock:                       # dong 585: atomic
            for key, value in patch.items():
                setattr(self.live_cfg, key, value)# dong 587: hot-reload
\end{lstlisting}
\end{code}

\begin{say}
``\textbf{Dòng 515}: \texttt{\_\_init\_\_} nhận camera index và preset name.
\texttt{stop\_event} là threading.Event --- cờ hiệu để dừng pipeline.
\texttt{cfg\_lock} là Lock để protect live config khi hot-reload.

\textbf{Dòng 571}: \texttt{start()} tạo thread daemon, gọi \texttt{self.\_run()}.
Daemon thread: tự chết khi main process chết --- không leak process.

\textbf{Dòng 575}: \texttt{stop()} --- hai bước quan trọng:
Dòng 577: \texttt{stop\_event.set()} --- set cờ.
Thread \texttt{\_run()} kiểm tra cờ này ở cuối mỗi vòng loop và tự thoát.
Dòng 578: \texttt{join(timeout=10.0)} --- chờ tối đa 10 giây để thread thoát gracefully.
Timeout để không block mãi mãi nếu thread bị stuck.

\textbf{Dòng 583--587}: \texttt{patch\_config} --- hot-reload config không restart pipeline.
Client frontend gửi PATCH /pipeline/config.
\texttt{with self.cfg\_lock}: đảm bảo atomic --- thread pipeline không đọc config lúc đang write.
\texttt{setattr}: thay đổi key-value bất kỳ trong \texttt{live\_cfg}.
Ví dụ: đổi từ balanced sang strong, hoặc đổi detect\_every từ 2 sang 1.
Pipeline đang chạy nhận config mới ngay frame tiếp theo --- không cần restart.''
\end{say}

\subsection{4:30--6:00 --- \_run: processing loop (dòng 636--750)}

\begin{doact}
Nhấn \texttt{Ctrl+G}, gõ \textbf{636}, Enter.
\end{doact}

\begin{code}{app\_api.py dòng 636--752 --- \_run(): main processing loop}
\begin{lstlisting}[language=Python, firstnumber=636]
def _run(self) -> None:
    """Main pipeline loop running in background thread."""
    cam_idx  = self.camera_index
    cap      = cv2.VideoCapture(cam_idx)     # dong 639

    # --- Khoi dong 3 module tu cac file khac ---
    detector   = FaceDetector(               # dong 674: tu Khoa
        model_path=MODEL_PATH,
        conf_threshold=cfg.conf_threshold,
        iou_threshold=cfg.iou_threshold,
        imgsz=cfg.imgsz)
    db         = FamilyDatabase(             # dong 688: tu Khoa
        db_path=str(_P("data/family_members.json")))
    recognizer = FaceRecognizer(             # dong 692: tu Khoa
        model=cfg.recognition_model)

    frame_reader = FrameReader(cap)          # dong 711: tu run_phone_cam.py
    frame_reader.start()                     # dong 712: thread start

    frame_idx = 0

    while not self.stop_event.is_set():      # dong 738: loop cho den khi stop
        ok, frame = frame_reader.read()      # dong 739: non-blocking read
        if not ok:
            time.sleep(0.04)
            continue
        frame_idx += 1

        if frame_idx % cfg.detect_every == 0:  # dong 750: YOLO every N frames
            boxes_scores = detector.detect_faces_with_scores(frame)
\end{lstlisting}
\end{code}

\begin{say}
``\textbf{Dòng 639}: \texttt{cv2.VideoCapture(cam\_idx)} mở kết nối camera.

\textbf{Dòng 674}: \texttt{FaceDetector} của Khoa. Khởi động một lần, tái dùng mãi mãi.
\textbf{Dòng 688}: \texttt{FamilyDatabase} của Khoa. Load JSON từ disk vào RAM.
\textbf{Dòng 692}: \texttt{FaceRecognizer} của Khoa. Load InsightFace buffalo\_l model.

Ba module này \textbf{em không viết} --- em chỉ orchestrate chúng ở đây.
Đây là Separation of Concerns: Khoa viết class riêng, em gọi từ pipeline.

\textbf{Dòng 711--712}: \texttt{FrameReader(cap)} và \texttt{start()} --- khởi động thread đọc frame.
\texttt{FrameReader} em giải thích ngay sau.

\textbf{Dòng 738}: \texttt{while not self.stop\_event.is\_set()} --- loop tới khi \texttt{stop()} gọi.

\textbf{Dòng 750}: \texttt{frame\_idx \% cfg.detect\_every == 0} --- chỉ chạy YOLO mỗi N frame.
\texttt{detect\_every=2} (preset balanced): YOLO chạy 15fps, display vẫn 30fps.
Tiết kiệm 50\% CPU cho YOLO --- box cũ vẫn apply cho frame bỏ qua.''
\end{say}

% ============================================================
\section{6:00--7:30 --- RUN\_PHONE\_CAM.PY --- FRAMEREADER}
% ============================================================

\begin{timing}{6:00 -- 7:30}
Chuyển sang tab \fpath{src/run\_phone\_cam.py}.
\end{timing}

\begin{doact}
Click tab \fpath{run\_phone\_cam.py}.
Nhấn \texttt{Ctrl+G}, gõ \textbf{327}, Enter.
\end{doact}

\begin{say}
``Điểm thứ ba: \texttt{FrameReader} class --- tại sao cần một thread riêng để đọc camera?''
\end{say}

\begin{code}{run\_phone\_cam.py dòng 327--382 --- FrameReader: non-blocking camera reader}
\begin{lstlisting}[language=Python, firstnumber=327]
class FrameReader(threading.Thread):
    """
    Dedicated thread to drain frames from VideoCapture.
    Keeps queue size = 1 so pipeline always gets THE LATEST frame,
    not a stale one from 500ms ago.
    """

    def __init__(self, cap: cv2.VideoCapture) -> None:
        super().__init__(daemon=True)              # dong 331
        self.cap     = cap
        self.queue: Queue = Queue(maxsize=1)       # dong 332: chi 1 slot
        self.running = True

    def run(self) -> None:                        # dong 338
        while self.running:
            ok, frame = self.cap.read()           # dong 341: blocking
            if not ok:
                time.sleep(0.02)
                continue
            # Drop-stale pattern: loai bo frame cu neu queue day
            if self.queue.full():                 # dong 365
                try:
                    self.queue.get_nowait()       # dong 366: bỏ frame cu
                except Empty:
                    pass
            try:
                self.queue.put_nowait(frame)      # dong 369: push moi nhat
            except Full:
                pass

    def read(self, timeout: float = 0.6):         # dong 376
        try:
            return True, self.queue.get(          # dong 378
                timeout=timeout)
        except Empty:
            return False, None                    # dong 380

    def stop(self) -> None:                       # dong 382
        self.running = False
\end{lstlisting}
\end{code}

\begin{say}
``\textbf{Vấn đề nếu không có FrameReader:}
\texttt{cv2.VideoCapture.read()} là blocking call --- nó block thread cho tới khi camera có frame.
Nếu pipeline đang xử lý (YOLO + anonymize) mất 80ms,
\texttt{cap.read()} trong cùng thread phải chờ.
Trong 80ms đó camera có 2--3 frame mới --- tất cả bị drop trong OpenCV buffer.
Khi \texttt{cap.read()} tiếp theo: trả về frame từ 80ms trước --- \textbf{stale}.
Stream nhìn như lag.

\textbf{Giải pháp --- FrameReader:}
\textbf{Dòng 327}: inherit \texttt{threading.Thread}.
\textbf{Dòng 331}: \texttt{daemon=True} --- tự chết khi process chết.
\textbf{Dòng 332}: \texttt{Queue(maxsize=1)} --- chỉ 1 slot.
Không phải 30 slot (framerate buffer) --- chỉ 1.

\textbf{Dòng 341}: \texttt{cap.read()} chạy trong thread riêng.
Nó block --- nhưng không ảnh hưởng pipeline thread.

\textbf{Dòng 365--369}: đây là \textbf{drop-stale pattern}:
Nếu queue đã đầy (có frame chưa được lấy):
Bỏ frame cũ đó đi (dòng 366).
Đẩy frame mới vào (dòng 369).
Kết quả: queue luôn chứa frame \textbf{mới nhất}, không phải frame cũ.

\textbf{Dòng 376}: pipeline gọi \texttt{frame\_reader.read(timeout=0.6)}.
Nếu trong 0.6s không có frame (camera disconnect): trả False.
Không block pipeline mãi mãi.''
\end{say}

% ============================================================
\section{7:30--9:30 --- STREAMING ENDPOINTS}
% ============================================================

\begin{timing}{7:30 -- 9:30}
Quay về tab \fpath{app\_api.py}.
\end{timing}

\begin{doact}
Nhấn \texttt{Ctrl+G}, gõ \textbf{1540}, Enter.
\end{doact}

\begin{say}
``Điểm thứ tư: làm sao browser nhận video real-time từ server?
Ba protocol: MJPEG, SSE, WebSocket.''
\end{say}

\subsection{7:30--8:15 --- MJPEG Stream (dòng 1540--1572)}

\begin{code}{app\_api.py dòng 1540--1572 --- GET /stream/mjpeg: MJPEG multipart stream}
\begin{lstlisting}[language=Python, firstnumber=1540]
@app.get("/stream/mjpeg")
async def stream_mjpeg(
    session_id: Optional[str] = None,
    ...
) -> StreamingResponse:

    async def generate():
        while True:
            frame_data = await asyncio.get_event_loop(
            ).run_in_executor(None,
                session.frame_buffer.get, 0.5)  # dong 1553: get frame

            if frame_data is None:
                await asyncio.sleep(0.04)        # dong 1555: 25fps cap
                continue

            # MJPEG boundary format: RFC 2046 multipart
            yield (b"--frame\r\n"                # dong 1560: boundary
                   b"Content-Type: image/jpeg\r\n\r\n"
                   + frame_data + b"\r\n")
            await asyncio.sleep(0.04)            # dong 1563: 25fps

    return StreamingResponse(
        generate(),
        media_type="multipart/x-mixed-replace"  # dong 1567
                   ";boundary=frame",
    )
\end{lstlisting}
\end{code}

\begin{say}
``MJPEG là cách đơn giản nhất để stream video qua HTTP.

\textbf{Dòng 1560}: \texttt{--frame\textbackslash r\textbackslash n} là boundary separator.
Server gửi sequence: [boundary][header][jpeg bytes][boundary][header][jpeg bytes]...
Browser <img> tag decode từng JPEG liên tiếp --- thành video.

\textbf{Dòng 1555, 1563}: \texttt{asyncio.sleep(0.04)} = 25fps cap.
Tại sao 25fps chứ không 30? Camera 30fps, pipeline xử lý $\approx$25fps.
Cap ở 25 tránh buffer overflow.

\textbf{Tại sao không WebSocket luôn?}
MJPEG với \texttt{<img>} tag không cần JS --- works với \texttt{<img src="/stream/mjpeg">}.
Lowest latency setup. WebSocket cần JS to rebuild frames.''
\end{say}

\subsection{8:15--9:30 --- SSE và WebSocket (dòng 1581, 1620)}

\begin{doact}
Nhấn \texttt{Ctrl+G}, gõ \textbf{1581}, Enter. Đọc 5 giây. Nhấn \texttt{Ctrl+G}, gõ \textbf{1620}, Enter.
\end{doact}

\begin{say}
``\textbf{Dòng 1581}: \texttt{GET /stream/sse} --- Server-Sent Events.
Thay vì stream JPEG bytes, stream JSON metadata: bao nhiêu khuôn mặt, tên ai đã nhận ra, thống kê.
Frontend dùng để update số liệu dashboard mà không cần reload.

\textbf{Dòng 1620}: \texttt{@app.websocket("/stream/ws")} --- WebSocket.
Dùng khi client muốn bidirectional: subscribe stream và đồng thời gửi lệnh.
Ví dụ: đang stream, client gửi \texttt{\{"action":"unlock","member\_id":"dat"\}}.
Server nhận, unlock track đó ngay, stream tiếp.

Ba protocol cover mọi use case: MJPEG cho display đơn giản, SSE cho metadata, WebSocket cho control.''
\end{say}

% ============================================================
\section{9:30--11:00 --- CRUD ENDPOINTS (bảng tổng hợp)}
% ============================================================

\begin{timing}{9:30 -- 11:00}
Mở \texttt{localhost:8000/docs} trên browser --- Swagger UI.
\end{timing}

\begin{say}
``Điểm thứ năm: tổ chức endpoint.

Em gom 40+ endpoint thành groups rõ ràng.
Mở Swagger để thấy full list.''
\end{say}

\begin{doact}
Chuyển sang browser tab Swagger UI (\texttt{localhost:8000/docs}).
Scroll qua các group endpoint.
\end{doact}

\begin{tcolorbox}[colback=MidnightBlue!4, colframe=MidnightBlue!60, breakable,
  title={\faTable\ \textbf{Tổng hợp endpoints chính (trích dẫn dòng)}}]
\renewcommand{\arraystretch}{1.4}
\begin{tabularx}{\textwidth}{lllX}
\toprule
Method & Endpoint & Dòng (app\_api.py) & Mục đích \\
\midrule
GET  & /health          & 968  & Kiểm tra server alive \\
GET  & /presets         & 988  & Trả PRESETS dict của Đạt \\
POST & /anonymize       & 1010 & Ẩn danh 1 ảnh tĩnh (no pipeline) \\
POST & /lock            & 1035 & Lock track khỏi whitelist \\
POST & /unlock          & 1074 & Unlock track về whitelist \\
POST & /recognize       & 1102 & Nhận dạng 1 ảnh tĩnh \\
\midrule
GET  & /members         & 1261 & Danh sách thành viên đã đăng ký \\
POST & /members         & 1305 & Đăng ký thành viên mới \\
PATCH & /members/\{id\} & 1375 & Cập nhật thông tin / thêm embedding \\
DELETE & /members/\{id\} & 1408 & Xoá thành viên \\
\midrule
POST & /pipeline/start  & 1432 & Khởi động PipelineSession (dòng 512) \\
DELETE & /pipeline/stop & 1458 & Dừng pipeline, join thread (dòng 575) \\
GET  & /pipeline/status & 1474 & Trạng thái: alive/stopped, frame count \\
PATCH & /pipeline/config & 1488 & Hot-reload config (gọi patch\_config dòng 583) \\
\midrule
GET  & /stream/mjpeg    & 1540 & MJPEG stream (25fps) \\
GET  & /stream/sse      & 1581 & SSE metadata stream \\
WS   & /stream/ws       & 1620 & WebSocket bidirectional \\
\midrule
POST & /admin/setup     & 1707 & Tạo admin lần đầu (delegate sang Đạt auth) \\
POST & /admin/verify    & 1725 & Login admin, issue JWT token \\
\midrule
GET  & /clips           & 1763 & Danh sách clips đã ghi \\
GET  & /clips/\{id\}/video & 1792 & Download clip (raw, không decrypt) \\
POST & /clips/\{id\}/decrypt & 1835 & Decrypt clip với password \\
DELETE & /clips/\{id\} & 1876 & Xoá một clip \\
DELETE & /clips          & 1884 & Xoá tất cả clips \\
\bottomrule
\end{tabularx}
\end{tcolorbox}

\begin{say}
``Mỗi endpoint em sẽ chỉ một ví dụ quan trọng: \textbf{PATCH /pipeline/config}.''
\end{say}

\begin{doact}
Quay về VS Code, \texttt{Ctrl+G}, gõ \textbf{1488}, Enter.
\end{doact}

\begin{code}{app\_api.py dòng 1488--1504 --- PATCH /pipeline/config: hot-reload runtime}
\begin{lstlisting}[language=Python, firstnumber=1488]
@app.patch("/pipeline/config",
           dependencies=[Depends(_require_admin_auth)])  # dong 1489: auth
async def patch_pipeline_config(
    patch: Dict[str, Any],
    session_id: Optional[str] = None,
) -> dict:
    session = _get_session(session_id)
    if session is None:
        raise HTTPException(404, "No active pipeline")

    allowed = {"preset_name", "detect_every",    # dong 1498: whitelist
               "blur_scale", "pixel_block",
               "head_ratio", "anonymize"}
    for key in patch:
        if key not in allowed:
            raise HTTPException(400, f"Unknown key: {key}")
    session.patch_config(patch)                  # dong 1504: atomic update
    return {"status": "ok", "applied": patch}
\end{lstlisting}
\end{code}

\begin{say}
``\textbf{Dòng 1489}: \texttt{Depends(\_require\_admin\_auth)} --- admin only.
\textbf{Dòng 1498}: whitelist các keys được phép patch.
Tránh client truyền arbitrary attributes vào config object --- injection prevention.
\textbf{Dòng 1504}: gọi \texttt{session.patch\_config(patch)} --- atomic update dòng 583.
Response trả lại chính xác những gì đã apply --- cho client confirm.''
\end{say}

% ============================================================
\section{11:00--12:30 --- DEMO LIVE --- FRONTEND TOUR}
% ============================================================

\begin{timing}{11:00 -- 12:30}
Chứng minh toàn bộ stack hoạt động end-to-end.
\end{timing}

\begin{doact}
Mở browser tab 1: \texttt{http://localhost:8000/admin/dashboard.html}.
\end{doact}

\begin{say}
``Demo nhanh 3 trang.''
\end{say}

\begin{doact}
\textbf{Demo 1 --- Dashboard (dòng 1432, 1488, 1540):}\\
1. Click ``Start Pipeline'' $\to$ POST /pipeline/start dòng 1432 $\to$ PipelineSession.start() dòng 571.\\
2. Video stream hiện ra $\to$ MJPEG từ /stream/mjpeg dòng 1540.\\
3. Dropdown đổi preset ``strong'' $\to$ PATCH /pipeline/config dòng 1488 $\to$ patch\_config dòng 583.\\
4. Thấy blur mode thay đổi ngay không restart.\\
Nói: ``Hot-reload, không cần restart pipeline.''
\end{doact}

\begin{doact}
\textbf{Demo 2 --- Face Registry (dòng 1305, 1261):}\\
Chuyển sang tab face-registry.html.\\
Thấy danh sách thành viên $\to$ GET /members dòng 1261.\\
Nói: ``Đây là database của Khoa's family\_database.py được expose qua endpoint của em.''
\end{doact}

\begin{doact}
\textbf{Demo 3 --- Secure Archive (dòng 1763, 1835):}\\
Chuyển sang tab secure-archive.html.\\
Danh sách encrypted clips $\to$ GET /clips dòng 1763.\\
Bấm decrypt trên một clip, nhập password $\to$ POST /clips/\{id\}/decrypt dòng 1835.\\
Nói: ``Clip được encrypt bởi Đạt's clip\_recorder.py AES-256-GCM. Endpoint này gọi decrypt\_clip() của Đạt.''
\end{doact}

% ============================================================
\section{12:30--15:00 --- KẾT LUẬN TOÀN ĐỀ ÁN}
% ============================================================

\begin{timing}{12:30 -- 15:00}
QUAN TRỌNG NHẤT. Nhìn camera. Nói chậm, rõ. Đây là phần kết thúc toàn bộ đề án.
\end{timing}

\begin{say}
``Em là người trình bày cuối --- vậy em xin kết luận toàn bộ đề án.

\textbf{Ba thành viên, ba lớp trong pipeline:}

Đạt viết \textbf{lớp Privacy và Security}:
\fpath{privacy.py} 8 thuật toán ẩn danh,
\fpath{anonymizer\_backend.py} 4 preset cấu hình,
\fpath{clip\_recorder.py} AES-256-GCM encrypted storage,
\fpath{admin\_auth.py} PBKDF2 390.000 iterations + timing-safe compare.
Tổng 749 dòng.

Khoa viết \textbf{lớp Detection và Recognition}:
\fpath{detector.py} YOLOv11n-face,
\fpath{face\_recognition.py} InsightFace buffalo\_l embedding 512-dim,
\fpath{family\_database.py} vectorized cosine matching,
\fpath{bbox\_smoother.py} EMA tracking.
Tổng 759 dòng.

Em viết \textbf{lớp Orchestration và API}:
\fpath{app\_api.py} 40+ endpoints, PipelineSession, JWT middleware,
3 streaming protocol, hot-reload config,
\fpath{run\_phone\_cam.py} FrameReader non-blocking drop-stale pattern.
Tổng 3207 dòng.

\textbf{Tổng toàn đề án: $749 + 759 + 3207 = 4715$ dòng code.}

\pause\ \camlook\

\textbf{Ba quyết định thiết kế quan trọng nhất:}

Thứ nhất: \textbf{Privacy by Design}.
Không có option ``nhận diện trước, quyết định blur sau''.
Mọi khuôn mặt detect được đều bị xử lý ngay.
Whitelist (thành viên gia đình) là exception, không phải default.

Thứ hai: \textbf{Separation of Concerns}.
Privacy logic ở \fpath{privacy.py}. Detection ở \fpath{detector.py}.
Recognition ở \fpath{family\_database.py}. Orchestration ở \fpath{app\_api.py}.
Mỗi file có thể test, thay thế, nâng cấp độc lập.
Ví dụ: thay YOLOv11 bằng YOLOv12 $\to$ chỉ sửa \fpath{detector.py}, pipeline không đổi.

Thứ ba: \textbf{Security theo chuẩn OWASP}.
Không hardcode credentials. PBKDF2 390.000 iterations chống brute-force.
AES-256-GCM với salt ngẫu nhiên mỗi clip.
Timing-safe compare chống timing attack.
JWT check trên mọi admin endpoint.

\pause\ \camlook\

\textbf{Hạn chế và hướng phát triển:}

Hạn chế: buffalo\_l model 500MB --- cold start chậm.
InsightFace chạy CPU: 50--80ms/inference, nếu nhiều người trong frame đồng thời có thể lag.
Chưa có caching embedding thực sự nhanh (FAISS, Annoy).

Hướng phát triển:
Thay InsightFace bằng model nhỏ hơn cho edge device.
Thêm FAISS index cho database lớn hơn.
Certificate-based auth thay JWT base64 đơn giản.

\pause\ \camlook\

Đây là toàn bộ đề án XLA Face Privacy System.
Cảm ơn thầy cô đã theo dõi.''
\end{say}

% ============================================================
\section{PHỤ LỤC --- CÂU HỎI DỰ KIẾN CHO KHÁNH}
% ============================================================

\begin{tcolorbox}[colback=gray!5, colframe=MidnightBlue!70, breakable,
  title={\textbf{Q: Tại sao FastAPI chứ không Flask hay Django?}}]
``FastAPI: async-native --- quan trọng cho streaming (MJPEG, SSE, WebSocket) cùng lúc nhiều client.
Flask: sync mặc định, streaming phải dùng monkey-patch.
Django: too heavy cho task này, ORM không cần thiết.
FastAPI auto-generate Swagger UI từ Python type hints --- zero effort documentation.''
\end{tcolorbox}

\begin{tcolorbox}[colback=gray!5, colframe=MidnightBlue!70, breakable,
  title={\textbf{Q: JWT token base64 không phải signed --- có an toàn không?}}]
``Đúng --- đây là điểm hạn chế của implementation hiện tại.
Token chỉ base64 encode, không sign bằng secret key.
Ai decode được token có thể forge token hết hạn $\to$ đặt exp tương lai.
Trong phạm vi đề án nội bộ (LAN, không expose internet) chấp nhận được.
Production: dùng PyJWT với HS256 secret, thêm refresh token.''
\end{tcolorbox}

\begin{tcolorbox}[colback=gray!5, colframe=MidnightBlue!70, breakable,
  title={\textbf{Q: detect\_every=2 có làm miss khuôn mặt không?}}]
``Không miss --- YOLO chạy lần 2, box kết quả apply cho cả frame 1 và frame 2.
Frame lẻ dùng box từ lần detect trước kết hợp BBoxSmoother của Khoa.
Người di chuyển ở 30fps: 2 frame = 66ms.
Ở 66ms người di chuyển $\approx$3--5cm --- box prediction vẫn valid.
Nếu cần accuracy cao hơn: set detect\_every=1 qua PATCH /pipeline/config.''
\end{tcolorbox}

\begin{tcolorbox}[colback=gray!5, colframe=MidnightBlue!70, breakable,
  title={\textbf{Q: FrameReader Queue(maxsize=1) có bị mất frame clip không?}}]
``FrameReader phục vụ \textbf{live display} --- chỉ cần frame mới nhất.
ClipRecorder (của Đạt, dòng 248) có queue riêng, maxsize=300 (10s x 30fps).
Two separate queues: FrameReader queue for display (latest), ClipRecorder queue for recording (all frames).
Không dùng chung queue $\to$ không conflict.''
\end{tcolorbox}

\begin{tcolorbox}[colback=gray!5, colframe=MidnightBlue!70, breakable,
  title={\textbf{Q: Làm sao chạy thực tế? Cần GPU không?}}]
``Không cần GPU.
YOLOv11n-face trên CPU Intel i5 gen 10: $\approx$25--30ms/frame.
InsightFace buffalo\_l CPU: $\approx$50--80ms/frame.
Tổng pipeline latency: 80--120ms end-to-end.
Display lag $\approx$0.1s --- mắt thường không nhận ra.
Máy yếu hơn: tăng detect\_every lên 3--4, giảm imgsz xuống 416.''
\end{tcolorbox}

\end{document}
